% ----------------------------------------------------------------------
%  1.2.5. Матрица ортогонального преобразования; 1.2.6. Преобразование
%  индексов; 1.2.7. Формулы структурной кристаллографии; 1.2.8. Ячейки
%  Вигнера–Зейтца
% ----------------------------------------------------------------------
\subsection{Матрица ортогонального преобразования}\label{ssec:orth}

Повороты (и вообще поворотные элементы симметрии) описывают
\term{ортогональными матрицами}~— они сохраняют длины векторов и углы
между ними. Для ортогональной матрицы $A=\{a_{ij}\}$
\begin{equation}\label{eq:orth}
\sum_{i=1}^3a_{ik}a_{il}=\delta_{kl},\qquad A^{-1}=\tilde A,\qquad\det A=\pm1
\end{equation}
($\tilde A$~— транспонированная матрица). При $\det A=+1$ это
\emph{собственный} поворот, при $\det A=-1$~— \emph{несобственный}
(поворот с отражением или инверсией).

\paragraph{Ось и угол поворота по матрице}
Собственные значения вещественной ортогональной матрицы $3\times3$
(собственного поворота)~— $1$, $e^{-i\alpha}$, $e^{+i\alpha}$. Собственный вектор
$\vb u$ с собственным значением $1$ не меняется при повороте~— это
\emph{ось}, а фаза $\alpha$~— \emph{угол} поворота. Их находят прямо по
элементам матрицы:
\begin{keyeq}\label{eq:axis-angle}
\begin{gathered}
\operatorname{tr}A=a_{11}+a_{22}+a_{33}=1+2\cos\alpha,\\
u_1:u_2:u_3=(a_{32}-a_{23}):(a_{13}-a_{31}):(a_{21}-a_{12}) .
\end{gathered}
\end{keyeq}
След матрицы (по-английски trace, по-немецки Spur, $\operatorname{Sp}A$)~—
инвариант: он не зависит от системы координат.

\begin{primer}[Поворот вокруг {[111]}]
Матрица циклической перестановки осей $x\to y\to z\to x$:
\[
A=\begin{pmatrix}0&0&1\\1&0&0\\0&1&0\end{pmatrix}.
\]
$\operatorname{tr}A=0=1+2\cos\alpha$, откуда $\cos\alpha=-\frac12$, $\alpha=120^\circ$.
Ось: $(a_{32}-a_{23}):(a_{13}-a_{31}):(a_{21}-a_{12})=(1-0):(1-0):(1-0)$, т.\,е.
$[111]$. Это поворотная ось 3-го порядка куба (п.~1.3.1).
\end{primer}

Обратно, матрица поворота на угол $\varphi$ вокруг оси с направляющими
косинусами $l_1$, $l_2$, $l_3$\ispr{В оригинале (формула (1.60)) в элементе
$R_{12}$ стоит $-l_3\cos\varphi$; правильно $-l_3\sin\varphi$.}:
\begin{equation}\label{eq:R-axis}
R=\begin{pmatrix}
l_1^2(1-c)+c & l_1l_2(1-c)-l_3s & l_1l_3(1-c)+l_2s\\
l_1l_2(1-c)+l_3s & l_2^2(1-c)+c & l_2l_3(1-c)-l_1s\\
l_1l_3(1-c)-l_2s & l_2l_3(1-c)+l_1s & l_3^2(1-c)+c
\end{pmatrix},\quad c=\cos\varphi,\ s=\sin\varphi
\end{equation}

\paragraph{Углы Эйлера}
Если поворот задан углами Эйлера $\varphi$, $\theta$, $\psi$ (п.~\ref{ssec:projections}),
каждый из трёх поворотов в своей системе координат выглядит просто
(обозначим $c_x\equiv\cos x$, $s_x\equiv\sin x$):
\[
R_z(\varphi)=\begin{pmatrix}c_\varphi&-s_\varphi&0\\s_\varphi&c_\varphi&0\\0&0&1\end{pmatrix},\quad
R_{x'}(\theta)=\begin{pmatrix}1&0&0\\0&c_\theta&-s_\theta\\0&s_\theta&c_\theta\end{pmatrix},\quad
R_{z'}(\psi)=\begin{pmatrix}c_\psi&-s_\psi&0\\s_\psi&c_\psi&0\\0&0&1\end{pmatrix}
\]
В неподвижной системе координат матрицы перемножают в порядке,
\emph{обратном} порядку поворотов: $R(\varphi,\theta,\psi)=R_z(\varphi)R_{x'}(\theta)R_{z'}(\psi)$,
\begin{equation}\label{eq:euler}
R=\begin{pmatrix}
c_\varphi c_\psi-c_\theta s_\varphi s_\psi & -c_\varphi s_\psi-c_\theta s_\varphi c_\psi & s_\theta s_\varphi\\
s_\varphi c_\psi+c_\theta c_\varphi s_\psi & -s_\varphi s_\psi+c_\theta c_\varphi c_\psi & -s_\theta c_\varphi\\
s_\theta s_\psi & s_\theta c_\psi & c_\theta
\end{pmatrix}
\end{equation}
($0\le\varphi,\psi<2\pi$, $0\le\theta\le\pi$)\ispr{В оригинале (формула (1.61)) в
элементе $R_{11}$ пропущен множитель $\cos\theta$ при
$\sin\varphi\sin\psi$.}.
Углы Эйлера популярны в физике твёрдого тела (и при описании текстур,
п.~1.6.2); кристаллографа же чаще интересуют индексы оси и угол поворота.

\paragraph{Активная и пассивная трактовки}
Вектор $\vb r$ при повороте переходит в $\vb r'=A\vb r$. Эту запись можно
понимать двояко (рис.~\ref{fig:active-passive}): \emph{активно}~— вектор
повёрнут на угол $\alpha$, а система координат та же; \emph{пассивно}~—
вектор на месте, а система координат повёрнута на $\alpha$ в обратную
сторону. Удобное правило для составления матрицы: её \vazhno{столбцы~—
координаты векторов $[100]$, $[010]$, $[001]$ после поворота}.

\begin{figure}[ht]
\centering
\begin{subfigure}[b]{0.42\textwidth}\centering
\begin{tikzpicture}[scale=1.4]
  \draw[-{Stealth[length=2mm]}] (-0.2,0) -- (1.8,0) node[right, font=\small]{$y$};
  \draw[-{Stealth[length=2mm]}] (0,-0.2) -- (0,1.8) node[above, font=\small]{$z$};
  \draw[bvec, fmBlue] (0,0) -- (20:1.5) node[right, font=\small]{$\vb r$};
  \draw[bvec, fmRed] (0,0) -- (65:1.5) node[above right=-2pt, font=\small]{$\vb r'$};
  \draw[-{Stealth[length=1.5mm]}, fmGray] (20:0.7) arc (20:65:0.7);
  \node[font=\small, fmGray] at (42:0.9) {$\alpha$};
\end{tikzpicture}
\caption{}
\end{subfigure}\hfill
\begin{subfigure}[b]{0.42\textwidth}\centering
\begin{tikzpicture}[scale=1.4]
  \draw[-{Stealth[length=2mm]}] (-0.2,0) -- (1.8,0) node[right, font=\small]{$y$};
  \draw[-{Stealth[length=2mm]}] (0,-0.2) -- (0,1.8) node[above, font=\small]{$z$};
  \draw[-{Stealth[length=2mm]}, fmRed] (0,0) -- (-45:1.6) node[right, font=\small]{$y'$};
  \draw[-{Stealth[length=2mm]}, fmRed] (0,0) -- (45:1.6) node[above right=-2pt, font=\small]{$z'$};
  \draw[bvec, fmBlue] (0,0) -- (20:1.5) node[right, font=\small]{$\vb r$};
  \draw[-{Stealth[length=1.5mm]}, fmGray] (0:0.7) arc (0:-45:0.7);
  \node[font=\small, fmGray] at (-22:0.95) {$\alpha$};
\end{tikzpicture}
\caption{}
\end{subfigure}
\caption{Две трактовки поворота вокруг оси $x$: \textit{а}~— активная
(поворачивается вектор); \textit{б}~— пассивная (поворачиваются оси, в
обратную сторону): относительно новых осей вектор $\vb r$ имеет те же
координаты, что $\vb r'$ относительно старых}
\label{fig:active-passive}
\end{figure}

\paragraph{Смена системы координат}
Пусть в системе $xyz$ операция задана матрицей $A$: $\vb r_2=A\vb r_1$, а
переход к системе $x'y'z'$~— матрицей $C$: $\vb r'=C\vb r$. Тогда в новой
системе та же операция записывается матрицей
\begin{equation}\label{eq:similar}
A'=CAC^{-1}
\end{equation}
(подставьте $\vb r_{1,2}=C^{-1}\vb r'_{1,2}$ в $\vb r_2=A\vb r_1$). Такие матрицы
называют \emph{подобными}; все их инварианты, в том числе след, одинаковы:
$\operatorname{tr}A'=\operatorname{tr}A=1+2\cos\alpha$. Угол поворота от выбора осей,
естественно, не зависит.

\subsection{Преобразование индексов при смене системы координат}\label{ssec:index-transform}

Часто нужно перейти от одной элементарной ячейки к другой (например, от
кубической к примитивной, от гексагональной к ортогексагональной). Пусть
новые базисные векторы выражаются через старые матрицей $A$:
\begin{equation}\label{eq:basis-transform}
\vb a'_i=\sum_jA_{ij}\vb a_j\qquad(\text{кратко }\vb a'=A\vb a).
\end{equation}
Как изменятся индексы направлений и плоскостей?
\begin{itemize}
\item Вектор $\vb R=\sum m_i\vb a_i=\sum m'_i\vb a'_i$ не меняется, поэтому
индексы \emph{направлений} преобразуются «обратно» базису
(\emph{контравариантно}):
\begin{keyeq}\label{eq:m-transform}
\vb m'=(\tilde A)^{-1}\vb m .
\end{keyeq}
\item Скалярное произведение $hm_1+km_2+lm_3$ (номер плоскости в
семействе~\eqref{eq:plane-q}) тоже не меняется, откуда индексы
\emph{плоскостей} (компоненты векторов обратной решётки) преобразуются
так же, как базисные векторы (\emph{ковариантно}):
\begin{keyeq}\label{eq:h-transform}
\vb h'=A\vb h .
\end{keyeq}
\end{itemize}

\begin{primer}[Удвоенная ячейка]
Возьмём новые оси $\vb a'_1=\vb a_1+\vb a_2$, $\vb a'_2=-\vb a_1+\vb a_2$, $\vb a'_3=\vb a_3$
(ячейка, повёрнутая на $45^\circ$ и вдвое большая по площади основания):
$A=\begin{pmatrix}1&1&0\\\br1&1&0\\0&0&1\end{pmatrix}$. Плоскость $(110)$:
$\vb h'=A\vb h=(2,0,0)$, т.\,е. в новых осях это $(100)$ (точнее, каждая
вторая плоскость семейства $(200)$). Направление $[110]$:
$(\tilde A)^{-1}=\frac12\begin{pmatrix}1&1&0\\\br1&1&0\\0&0&2\end{pmatrix}$,
$\vb m'=(1,0,0)$~— это $[100]$, как и должно быть: $\vb a_1+\vb a_2=\vb a'_1$.
\end{primer}

\subsection{Основные формулы структурной кристаллографии}\label{ssec:formulas}

С метрической матрицей $G$ (п.~\ref{ssec:lattice}) все расчёты в любой
сингонии сводятся к нескольким формулам. Обозначим $\vb m=(m_1,m_2,m_3)^T$~—
индексы направления, $\vb h=(h,k,l)^T$~— индексы плоскости.

\begin{itemize}
\item \emph{Период идентичности}~— расстояние между ближайшими узлами
вдоль $[m_1m_2m_3]$:
\begin{equation}\label{eq:period}
I_{[m]}=\sqrt{\tilde{\vb m}G\vb m}\qquad\bigl(\text{куб: }I=a\sqrt{m_1^2+m_2^2+m_3^2}\bigr).
\end{equation}
\item \emph{Межплоскостное расстояние}: $1/d^2_{hkl}=|\vb H|^2=\tilde{\vb h}G^{-1}\vb h$.
Это выражение~— \term{квадратичная форма} сингонии:
\begin{keyeq}\label{eq:quadform}
\begin{gathered}
\frac1{d^2}=\frac{h^2+k^2+l^2}{a^2}\ \ \text{(куб)},\qquad
\frac1{d^2}=\frac{h^2}{a^2}+\frac{k^2}{b^2}+\frac{l^2}{c^2}\ \ \text{(ромб)},\\
\frac1{d^2}=\frac43\,\frac{h^2+hk+k^2}{a^2}+\frac{l^2}{c^2}\ \ \text{(гекс)}.
\end{gathered}
\end{keyeq}
Тетрагональная~— как ромбическая с $b=a$.
\item \emph{Угол между направлениями}
\begin{equation}\label{eq:angle-dirs}
\cos\varphi=\frac{\tilde{\vb m}^{(2)}G\vb m^{(1)}}{\sqrt{\tilde{\vb m}^{(1)}G\vb m^{(1)}}\sqrt{\tilde{\vb m}^{(2)}G\vb m^{(2)}}} ,
\end{equation}
\emph{угол между плоскостями}~— то же с $G^{-1}$ и индексами плоскостей;
\emph{угол между направлением и нормалью к плоскости}:
$\cos\varphi=\tilde{\vb h}\vb m\big/\bigl(\sqrt{\tilde{\vb h}G^{-1}\vb h}\sqrt{\tilde{\vb m}G\vb m}\bigr)$.
В кубической сингонии все три формулы сводятся к одной:
\begin{keyeq}\label{eq:cos-cubic}
\cos\varphi=\frac{h_1h_2+k_1k_2+l_1l_2}{\sqrt{h_1^2+k_1^2+l_1^2}\sqrt{h_2^2+k_2^2+l_2^2}} .
\end{keyeq}
\item \emph{Объём ячейки}: $V=(\vb a_1[\vb a_2\times\vb a_3])=\sqrt{\det G}$; для
обратной решётки $V^*=\sqrt{\det G^{-1}}=1/V$.
\item \emph{Переход от прямого базиса к обратному}: один и тот же вектор
имеет в прямом базисе компоненты $\vb m$, в обратном~— $\vb h=G\vb m$
(обратно $\vb m=G^{-1}\vb h$). В кубической сингонии $G=a^2I$, и направление
$[uvw]$ параллельно нормали к плоскости $(uvw)$.
\end{itemize}

\begin{primer}[Расчёты для меди и магния]
\textit{Медь} (ГЦК, $a=0{,}3615$~нм): $d_{111}=a/\sqrt3=0{,}2087$~нм,
$d_{200}=a/2=0{,}1808$~нм; угол между $(111)$ и $(100)$:
$\cos\varphi=1/\sqrt3$, $\varphi=54{,}74^\circ$; между $(111)$ и $(1\br11)$:
$\cos\varphi=1/3$, $\varphi=70{,}53^\circ$. Период вдоль $[110]$:
$a\sqrt2=0{,}511$~нм, но кратчайшее межатомное расстояние $a\sqrt2/2=0{,}256$~нм,
так как ГЦК-решётка центрирована.\\
\textit{Магний} (ГПУ, $a=0{,}321$, $c=0{,}521$~нм): $d_{10\br10}=a\sqrt3/2=0{,}278$~нм,
$d_{0002}=c/2=0{,}261$~нм.
\end{primer}

\subsection{Области Вороного, ячейки Вигнера--Зейтца, зоны Бриллюэна}\label{ssec:wigner}

Элементарную ячейку можно выбрать и иначе. Соединим нулевой узел с
ближайшими узлами и проведём через середины этих отрезков
перпендикулярные им плоскости. Они вырежут многогранник, параллельными
переносами которого заполняется всё пространство. Он содержит ровно
один узел (в центре), т.\,е. является примитивной ячейкой, но в отличие от
параллелепипеда обладает \emph{полной симметрией решётки}. В
кристаллографии это \term{область Вороного}, в физике твёрдого
тела~— \term{ячейка Вигнера--Зейтца}.

\begin{itemize}
\item \emph{ГЦК}: 12 ближайших соседей вдоль $\nfam{110}$; ячейка огранена 12
ромбами $\{110\}$~— \emph{ромбический додекаэдр} (рис.~\ref{fig:wigner}\,\textit{а}).
\item \emph{ОЦК}: 8 ближайших соседей вдоль $\nfam{111}$ и 6 следующих вдоль
$\nfam{100}$ (они лишь на 15\,\% дальше); ячейка огранена 8 правильными
шестиугольниками $\{111\}$ и 6 квадратами $\{100\}$~— 14-гранник,
\emph{усечённый октаэдр} (рис.~\ref{fig:wigner}\,\textit{б}).
\end{itemize}

Ячейка Вигнера--Зейтца, построенная в \emph{обратной} решётке,~—
\term{первая зона Бриллюэна} (гл.~3, п.~3.2.2). Поскольку ГЦК- и
ОЦК-решётки взаимно обратны, зона Бриллюэна ГЦК-кристалла~— усечённый
октаэдр, а ОЦК-кристалла~— ромбический додекаэдр.

\begin{figure}[ht]
\centering
\tikzset{bzfront/.style={atomLine, line width=0.7pt, line join=round},
  bzback/.style={atomLine!55, line width=0.4pt, dashed}}
\begin{subfigure}[b]{0.4\textwidth}\centering
\begin{tikzpicture}[scale=1.6]
  % сгенерировано tools/gen_bands.py — не править вручную
\fill[fmBlue!18] (-0.676,-0.539) -- (-0.469,-0.273) -- (-0.676,0.412) -- (-0.883,0.145) -- cycle;
\fill[fmBlue!14] (0.207,0.267) -- (-0.469,-0.273) -- (0.207,-0.684) -- (0.883,-0.145) -- cycle;
\fill[fmBlue!12] (0.207,-0.684) -- (-0.469,-0.273) -- (-0.676,-0.539) -- (0.000,-0.951) -- cycle;
\fill[fmBlue!20] (-0.676,0.412) -- (-0.469,-0.273) -- (0.207,0.267) -- (0.000,0.951) -- cycle;
\fill[fmBlue!8] (0.676,-0.412) -- (0.883,-0.145) -- (0.207,-0.684) -- (0.000,-0.951) -- cycle;
\fill[fmBlue!16] (0.207,0.267) -- (0.883,-0.145) -- (0.676,0.539) -- (0.000,0.951) -- cycle;
\draw[bzback] (0.469,0.273) -- (-0.207,0.684);
\draw[bzback] (0.469,0.273) -- (-0.207,-0.267);
\draw[bzback] (-0.883,0.145) -- (-0.207,-0.267);
\draw[bzback] (-0.883,0.145) -- (-0.207,0.684);
\draw[bzback] (0.469,0.273) -- (0.676,-0.412);
\draw[bzback] (0.469,0.273) -- (0.676,0.539);
\draw[bzfront] (0.883,-0.145) -- (0.676,0.539);
\draw[bzfront] (0.883,-0.145) -- (0.676,-0.412);
\draw[bzfront] (-0.469,-0.273) -- (-0.676,-0.539);
\draw[bzfront] (-0.469,-0.273) -- (-0.676,0.412);
\draw[bzfront] (-0.883,0.145) -- (-0.676,0.412);
\draw[bzfront] (-0.883,0.145) -- (-0.676,-0.539);
\draw[bzfront] (-0.469,-0.273) -- (0.207,0.267);
\draw[bzfront] (-0.469,-0.273) -- (0.207,-0.684);
\draw[bzfront] (0.883,-0.145) -- (0.207,-0.684);
\draw[bzfront] (0.883,-0.145) -- (0.207,0.267);
\draw[bzfront] (0.000,-0.951) -- (0.676,-0.412);
\draw[bzback] (0.000,-0.951) -- (-0.207,-0.267);
\draw[bzback] (0.000,0.951) -- (-0.207,0.684);
\draw[bzfront] (0.000,0.951) -- (0.676,0.539);
\draw[bzfront] (0.000,-0.951) -- (-0.676,-0.539);
\draw[bzfront] (0.000,-0.951) -- (0.207,-0.684);
\draw[bzfront] (0.000,0.951) -- (0.207,0.267);
\draw[bzfront] (0.000,0.951) -- (-0.676,0.412);
\coordinate (bzG) at (0.000,0.000);
\coordinate (bzH) at (-0.469,-0.273);
\coordinate (bzN) at (0.207,-0.209);
\coordinate (bzP) at (0.207,0.267);

\end{tikzpicture}
\caption{ГЦК}
\end{subfigure}\hspace{8mm}
\begin{subfigure}[b]{0.4\textwidth}\centering
\begin{tikzpicture}[scale=1.6]
  % сгенерировано tools/gen_bands.py — не править вручную
\fill[fmBlue!17] (-0.469,0.203) -- (-0.911,-0.200) -- (-0.469,-0.748) -- (-0.028,-0.345) -- cycle;
\fill[fmBlue!11] (0.883,-0.621) -- (1.118,-0.009) -- (0.883,0.330) -- (0.648,-0.281) -- cycle;
\fill[fmBlue!20] (0.441,0.879) -- (0.235,1.087) -- (-0.441,1.024) -- (-0.235,0.815) -- cycle;
\fill[fmBlue!15] (-0.911,-0.200) -- (-1.118,0.009) -- (-0.883,-0.330) -- (-0.441,-0.879) -- (-0.235,-1.087) -- (-0.469,-0.748) -- cycle;
\fill[fmBlue!21] (-0.235,0.815) -- (-0.441,1.024) -- (-0.883,0.621) -- (-1.118,0.009) -- (-0.911,-0.200) -- (-0.469,0.203) -- cycle;
\fill[fmBlue!10] (-0.469,-0.748) -- (-0.235,-1.087) -- (0.441,-1.024) -- (0.883,-0.621) -- (0.648,-0.281) -- (-0.028,-0.345) -- cycle;
\fill[fmBlue!17] (0.648,-0.281) -- (0.883,0.330) -- (0.441,0.879) -- (-0.235,0.815) -- (-0.469,0.203) -- (-0.028,-0.345) -- cycle;
\draw[bzfront] (-0.911,-0.200) -- (-0.469,0.203);
\draw[bzfront] (-0.911,-0.200) -- (-0.469,-0.748);
\draw[bzfront] (-0.469,-0.748) -- (-0.028,-0.345);
\draw[bzfront] (-0.469,0.203) -- (-0.028,-0.345);
\draw[bzback] (0.028,0.345) -- (0.469,-0.203);
\draw[bzback] (0.028,0.345) -- (0.469,0.748);
\draw[bzback] (0.469,0.748) -- (0.911,0.200);
\draw[bzback] (0.469,-0.203) -- (0.911,0.200);
\draw[bzfront] (1.118,-0.009) -- (0.883,-0.621);
\draw[bzfront] (1.118,-0.009) -- (0.883,0.330);
\draw[bzfront] (0.883,0.330) -- (0.648,-0.281);
\draw[bzfront] (0.883,-0.621) -- (0.648,-0.281);
\draw[bzback] (-0.648,0.281) -- (-0.883,0.621);
\draw[bzback] (-0.648,0.281) -- (-0.883,-0.330);
\draw[bzfront] (-0.883,-0.330) -- (-1.118,0.009);
\draw[bzfront] (-0.883,0.621) -- (-1.118,0.009);
\draw[bzfront] (0.235,1.087) -- (0.441,0.879);
\draw[bzfront] (0.235,1.087) -- (-0.441,1.024);
\draw[bzfront] (-0.441,1.024) -- (-0.235,0.815);
\draw[bzfront] (0.441,0.879) -- (-0.235,0.815);
\draw[bzback] (0.235,-0.815) -- (-0.441,-0.879);
\draw[bzback] (0.235,-0.815) -- (0.441,-1.024);
\draw[bzfront] (0.441,-1.024) -- (-0.235,-1.087);
\draw[bzfront] (-0.441,-0.879) -- (-0.235,-1.087);
\draw[bzback] (0.028,0.345) -- (-0.648,0.281);
\draw[bzback] (0.469,-0.203) -- (0.235,-0.815);
\draw[bzfront] (-0.883,-0.330) -- (-0.441,-0.879);
\draw[bzback] (0.469,0.748) -- (0.235,1.087);
\draw[bzfront] (-0.883,0.621) -- (-0.441,1.024);
\draw[bzback] (0.911,0.200) -- (1.118,-0.009);
\draw[bzfront] (0.441,-1.024) -- (0.883,-0.621);
\draw[bzfront] (0.441,0.879) -- (0.883,0.330);
\draw[bzfront] (-1.118,0.009) -- (-0.911,-0.200);
\draw[bzfront] (-0.235,-1.087) -- (-0.469,-0.748);
\draw[bzfront] (-0.235,0.815) -- (-0.469,0.203);
\draw[bzfront] (0.648,-0.281) -- (-0.028,-0.345);
\coordinate (bzG) at (0.000,0.000);
\coordinate (bzX) at (-0.469,-0.273);
\coordinate (bzL) at (0.207,0.267);
\coordinate (bzW) at (-0.028,-0.345);

\end{tikzpicture}
\caption{ОЦК}
\end{subfigure}
\caption{Ячейки Вигнера--Зейтца (расчёт; ось $z$~— вертикально):
\textit{а}~— ромбический додекаэдр ГЦК-решётки; \textit{б}~— усечённый
октаэдр ОЦК-решётки. Те же многогранники~— зоны Бриллюэна ОЦК- и
ГЦК-решёток соответственно}
\label{fig:wigner}
\end{figure}

\begin{itogi}
\begin{itemize}
\item Стереографическая проекция: направление~— точка ($\rho=r\operatorname{tg}\frac\theta2$),
плоскость~— дуга большого круга; гномостереографическая: плоскость~—
полюс. Углы измеряют по сеткам Вульфа и Болдырева.
\item Решётка: $\vb R=m_1\vb a_1+m_2\vb a_2+m_3\vb a_3$; 7 сингоний; метрическая
матрица $G=\{(\vb a_i\vb a_j)\}$.
\item Символы: узел $[\![mnp]\!]$, направление $[uvw]$, $\nfam{uvw}$; плоскость
$(hkl)$~— обратные отрезки; форма $\{hkl\}$; в гексагональной системе
$(hkil)$, $i=-(h+k)$.
\item Обратная решётка: $(\vb a_i\vb b_j)=\delta_{ij}$; $\vb H_{hkl}\perp(hkl)$,
$|\vb H_{hkl}|=1/d_{hkl}$; закон зон $hu+kv+lw=0$.
\item Поворот: $\operatorname{tr}A=1+2\cos\alpha$; индексы направлений
преобразуются контравариантно, плоскостей~— ковариантно.
\item $1/d^2=\tilde{\vb h}G^{-1}\vb h$; в кубе $d=a/\sqrt{h^2+k^2+l^2}$.
\item Ячейка Вигнера--Зейтца ГЦК~— ромбический додекаэдр, ОЦК~— усечённый
октаэдр.
\end{itemize}
\end{itogi}

\begin{proverka}
\item Чем стереографическая проекция плоскости отличается от
гномостереографической? Почему в кубическом кристалле полюс $(hkl)$
совпадает с выходом $[hkl]$?
\item Найдите расстояние от центра проекции $(001)$ до полюса $(011)$ при
$r=10$~см.
\item Определите индексы плоскости, отсекающей на осях отрезки $2a$,
$3b$, $-c$.
\item Переведите в четырёхиндексную запись направления $[210]$ и $[1\br10]$
гексагонального кристалла.
\item Найдите ось зоны плоскостей $(111)$ и $(1\br10)$.
\item Каков угол поворота, если след матрицы поворота равен $-1$?
\item Найдите межплоскостные расстояния $d_{110}$ и $d_{111}$ для
$\alpha$-железа ($a=0{,}2866$~нм).
\item Почему ячейку Вигнера--Зейтца ОЦК-решётки ограничивают не только
плоскости $\{111\}$, но и $\{100\}$?
\end{proverka}
